\chapter*{全书收束：从会运行到会负责}
\addcontentsline{toc}{chapter}{全书收束：从会运行到会负责}

这本书的主线并不是“学会更多工具”，而是从会运行代码，逐步走向能对一个科学计算项目负责。

Part I 建立了可复现科研计算的地基：目录、命令行、Git、脚本、数据入口和图表。Part II 让你学会把天文与物理数据读成有单位、有误差、有坐标、有采样边界的科学证据。Part III 把机器学习放进任务定义、baseline、评估和可信解释中。Part IV 通过案例把方法装配成小项目。Part V 讨论深度学习和现代 AI 工具，同时强调验证边界。Part VI 则把所有能力收束到 capstone 交付。

\section*{全书最终留下的一套项目包}

如果把整本书压缩成一个最终 portfolio，它不需要堆满所有中间草稿，但至少应该留下一个可以互相引用的 Capstone Evidence Package：

\begin{enumerate}
    \item Project summary：科学问题、范围和可回答性；
    \item Data Card / Dataset Contract：数据来源、字段、单位、误差、许可证和筛选；
    \item Evidence Records：关键操作、baseline、图表、诊断和复查记录；
    \item Model Experiment Record：模型选择、输入输出、训练验证和主要失败样本；
    \item Evaluation and error artifacts：指标、诊断图、误差分析和测试边界；
    \item Trust Statement：主要 claim、局限性、不确定度和 caveat；
    \item Reproducibility evidence：运行入口、环境、随机种子、Git 记录和复跑说明；
    \item Integrity and disclosure evidence：AI Usage Log、引用核查、公开边界和 human signoff。
\end{enumerate}

前六个 Part 的所有练习，最终都可以回到这一套 Evidence Package。它把“我做过很多 notebook”整理成“别人可以复查、评分、接手和维护的一份科学计算项目”。

\section*{最应该带走的习惯}

完成本书后，你不一定记得每一个命令、每一个公式或每一个模型参数。但下面这些习惯应当留下来：

\begin{itemize}
    \item 任何结果都要能追溯到数据、代码、环境和运行入口；
    \item 任何数值都要先问单位、误差和适用条件；
    \item 任何模型都要先有 baseline，再谈复杂度；
    \item 任何高分都要回到错误样本和科学代价中解释；
    \item 任何 AI 辅助输出都要经过验证、引用和责任边界检查；
    \item 任何项目都应为别人复查、接手和继续维护留下道路。
\end{itemize}

\section*{本科阶段最珍贵的能力}

本科阶段不需要把每一种现代 AI 方法都掌握到研究前沿。更珍贵的是形成一种可靠的工作方式：遇到新数据时知道先检查什么，遇到新模型时知道先比较什么，遇到漂亮结果时知道先怀疑什么，遇到 AI 助手时知道哪些地方可以借力、哪些地方必须自己负责。

这正是天文和物理训练与 AI 实战结合的意义。AI 可以放大你的能力，但也会放大你的疏忽。一个可靠的研究者，不是从不犯错，而是能设计工作流，让错误更早暴露、更容易定位、更不容易进入最终结论。

\section*{下一步}

如果你把本书作为课程使用，下一步可以选择一个 capstone 题目，从最小数据样本和最小 baseline 开始，逐周推进。如果你把本书作为自学材料，建议不要只阅读 PDF，而要真正运行 notebook、修改参数、制造错误、查看日志、提交版本，并把结果写成一份别人能读懂的小报告。

从这里开始，真正的学习不再是“看懂教材”，而是把教材中的工作流带进你自己的问题里。
